\section{Introduction}
\label{sec:introduction}

Access to legal information is a fundamental right in any state governed by
the rule of law. In Morocco, the legal corpus is vast and complex, comprising
hundreds of laws, codes, and decrees written primarily in Modern Standard
Arabic (MSA). For citizens, legal practitioners, and researchers, navigating
this corpus to find a precise answer to a specific question constitutes a
major challenge. General-purpose search engines (Google, Bing) are poorly
suited to this task: they neither understand the specificity of Arabic legal
vocabulary nor the hierarchical structure of legal codes (articles, chapters,
sections).

Recent advances in Natural Language Processing (NLP), particularly in Large
Language Models (LLMs) and Retrieval-Augmented Generation (RAG)
\cite{lewis2020retrieval}, offer an unprecedented opportunity to build
specialized Question Answering (Q\&A) systems. The RAG approach combines an
information-retrieval module with a text-generation module (LLM), producing
contextualized, sourced, and up-to-date responses without requiring expensive
fine-tuning. However, applying RAG to the Arabic legal domain presents unique
challenges: the scarcity of multilingual embedding models optimized for legal
vocabulary, the morphological complexity of Arabic \cite{habash2010arabic},
the presence of Darija (the Moroccan Arabic dialect) in user queries, the
need for absolute fidelity to legal texts, and the difficulty of processing
scanned historical legal documents.

Our system relies on a corpus of \textbf{2{,}357 legal articles} covering
seven areas of Moroccan law: criminal (822 articles), traffic (228), civil
(705), commercial (359), judicial organization (32), tax (211), and
violations. These articles are extracted from their source PDFs, cleaned,
normalized to standard Arabic, structured with hierarchical metadata, and
indexed in PostgreSQL + pgvector using the
\texttt{intfloat/multilingual-e5-large} embedding model (1024 dimensions).

Our contributions are fourfold:
\begin{enumerate}
  \item A \textbf{complete Arabic NLP processing chain} from PDF extraction
        (digital and scanned) to morphological normalization, cleaning, and
        structural analysis of Moroccan legal codes.
  \item A \textbf{hybrid search architecture} combining dense vector
        similarity (\texttt{BAAI/bge-m3}, 1024 dimensions) with BM25 lexical
        search, fused via Reciprocal Rank Fusion (RRF), indexed in
        PostgreSQL + pgvector with an HNSW index.
  \item A \textbf{systematic benchmark} of five multilingual embedding models
        on 26 Arabic legal articles with 65 query variants covering formal
        Arabic, Darija, French, English, and short/long queries.
  \item A \textbf{comparative benchmark} of three search strategies (dense,
        BM25, RRF) on 32 questions across 7 legal domains, evaluating both
        retrieval quality and generated-answer quality.
\end{enumerate}

The remainder of this paper is organized as follows. Section~\ref{sec:related}
presents related work. Section~\ref{sec:methodology} describes our
methodology. Section~\ref{sec:solution} details the proposed system
architecture. Section~\ref{sec:results} presents experimental results.
Section~\ref{sec:discussion} discusses implications and limitations.
Section~\ref{sec:conclusion} concludes the paper.
